% Conclusion section
\section{Conclusions}
\label{sec:conclusion}

This paper presented a GPU-accelerated numerical method for simulating droplet impact on curved surfaces using the Allen-Cahn phase-field lattice Boltzmann method with volume penalization and geometric gradient amplification. The key findings are:

\begin{enumerate}
    \item \textbf{Method validation}: The method successfully reproduces asymmetric spreading on cylindrical ridges with $k = 3.07$ at $R^* = 1.0$ (18.2\% error vs. experimental $k \approx 2.6$) and $k = 1.69$ at $R^* = 2.76$ (27.1\% error vs. experimental $k \approx 1.33$).
    
    \item \textbf{Grid convergence}: Resolution $N \geq 120$ provides stable, converged results when the geometric amplification parameter is properly calibrated. The amplification coefficient follows approximately $\text{amp} \propto N^{0.84}$.
    
    \item \textbf{Physical consistency}: The simulations show physically consistent behavior: $k$ decreases monotonically with increasing $R^*$ (converging to $k = 1$ for flat surfaces) and increases with Weber number following $k \propto \We^{0.80}$.
    
    \item \textbf{GPU efficiency}: The PyTorch CUDA implementation enables $150^3$ simulations at approximately 2.8 steps/sec on commodity hardware with 6 GB VRAM, making systematic parameter studies feasible.
    
    \item \textbf{Flat surface benchmark}: Symmetric spreading ($k = 1.0$) on flat surfaces confirms numerical accuracy and validates that asymmetric results on curved surfaces arise from physical geometric effects.
\end{enumerate}

\subsection{Limitations}

Several limitations should be addressed in future work:

\begin{enumerate}
    \item \textbf{Systematic discrepancy}: The simulated $k$ values consistently overestimate experimental measurements (18--27\% error). Possible causes include contact angle calibration, volume penalization smoothing scale, or Allen-Cahn mobility parameters. Further investigation is needed to reduce this gap.
    
    \item \textbf{Geometric amplification calibration}: The amplification parameter requires calibration for each grid resolution. A resolution-independent formulation or adaptive algorithm would improve the method's robustness.
    
    \item \textbf{Static contact angle limitation}: The Allen-Cahn sharpening term constrains achievable static contact angles to approximately $170^\circ$, whereas the target is $162^\circ$. The dynamic spreading predictions remain accurate despite this limitation.
    
    \item \textbf{Quasi-3D simulations}: The simulations use $N_z < N_x = N_y$ (typically $N_z \approx 0.7 N_x$) to reduce computational cost. True 3D effects on spreading asymmetry remain to be quantified.
    
    \item \textbf{No dynamic contact angle}: The current boundary condition uses only the equilibrium contact angle. Dynamic contact angle hysteresis and contact line velocity effects are not captured.
\end{enumerate}

\subsection{Future Directions}

Potential extensions of this work include:

\begin{itemize}
    \item Adaptive geometric amplification based on local interface curvature
    \item Investigation of alternative phase-field formulations (Cahn-Hilliard) for improved static contact angle prediction
    \item Extension to superhydrophobic surfaces with contact angle $> 150^\circ$
    \item Study of bouncing and partial rebound regimes on curved surfaces
    \item Multi-droplet interactions and coalescence dynamics
\end{itemize}

\section*{Acknowledgments}

This work was supported by the National Natural Science Foundation of China (grant number to be added).
The authors gratefully acknowledge the computational resources provided by Tencent Cloud.

\section*{Data Availability}

The simulation data and analysis scripts used in this work are available from the corresponding author upon reasonable request.